\begin{figure}[htbp]
\centering
\begin{tikzpicture}[causal,x=1cm,y=1cm]
 \node[vtx] (f12aX) at (0,0) {$X$};
 \node[vtx] (f12aY) at (2.6,0) {$Y$};
 \node[latent] (f12aZ) at (1.3,1.1) {$Z$};
 \draw[edd] (f12aZ)--(f12aX);
 \draw[edd] (f12aZ)--(f12aY);
 \node[vtx] (f12bX) at (5.0,0) {$X$};
 \node[vtx] (f12bY) at (7.0,0) {$Y$};
 \draw[ed] (f12bX)--(f12bY);
 \node[vtx] (f12cX) at (9.4,0) {$X$};
 \node[vtx] (f12cY) at (11.4,0) {$Y$};
 \draw[ed] (f12cY)--(f12cX);
\end{tikzpicture}
\caption{Reichenbach 的共同原因原理在统计性质与因果结构之间建立了联系。两个可观测变量 $X$ 和 $Y$ 之间的统计相依关系，表明它们受到某个变量 $Z$ 的共同影响；$Z$ 通常称为混杂因子（左图）。这里，$Z$ 也可以与 $X$ 或 $Y$ 中的一个重合，此时图形简化为中图或右图。该原理还指出，给定 $Z$ 时，$X$ 与 $Y$ 统计独立。图中的箭头表示直接因果关系，参见第 3 章和第 6 章。}
\label{fig:1.2}
\end{figure}
